\section{Relay Layer}
\label{sec:relay}

This section specifies TARDUS's optional relay layer (\Cref{sec:wallet} §6
referenced this layer as the canonical \emph{user-to-user delivery
mechanism} for off-chain coin transmission). The relay is a stateless,
TTL-bound, recipient-keyed mailbox: senders deposit ciphertexts addressed
to a recipient public key, and recipients poll. The relay never holds
key material and (under the v5.6 sealed-box payload format) never sees
plaintext coin material.

\subsection{Design Objectives}
\label{sec:relay-objectives}

\begin{description}[leftmargin=2em,style=nextline,itemsep=0pt]
  \item[O1 (recipient confidentiality)] A relay operator must not learn
    coin material (secrets, signatures, denomination) even on a full
    payload audit, except via decryption of the sealed-box payload —
    which requires the recipient's ed25519 secret key.
  \item[O2 (sender anonymity)] The relay must not authenticate the
    sender. The POST endpoint accepts unauthenticated requests; the
    only network-level identifier the relay sees is the connecting IP
    (mitigated at deployment by Tor/I2P or VPN tunnelling — out of
    scope for v5).
  \item[O3 (bounded resources)] Per-recipient and per-payload caps
    bound the relay's memory and disk usage; an oversized payload or
    a full inbox is rejected at the HTTP layer with 400 / 503.
  \item[O4 (eventual delivery)] Messages have a TTL (default 7 days);
    expired entries are pruned by a background task. The TTL forces
    recipients to poll within a bounded window; longer-term storage is
    the wallet's responsibility.
  \item[O5 (operator agnosticism)] The relay never inspects the
    payload's contents and never participates in the mint or refresh
    protocols. A relay outage delays delivery but cannot forge,
    double-spend, or otherwise compromise coin security.
\end{description}

\subsection{Wire Protocol}
\label{sec:relay-protocol}

The relay exposes three HTTP endpoints, all under the canonical path
\verb|/inbox/{recipient_pk_hex}|.

\paragraph{POST /inbox/\{recipient\_pk\_hex\}} — anonymous deposit.

\begin{description}[leftmargin=2em,itemsep=0pt]
  \item[Request]
    \verb|{ "payload_hex": "<bytes>", "ttl_secs": optional u64 }|
  \item[Response 200]
    \verb|{ "id": "<16-byte hex>", "payload_hex": "...", "received_at_unix_ms": u64, "expires_at_unix_ms": u64 }|
  \item[Response 400] payload exceeds \verb|--max-payload-bytes|, or
    \verb|recipient_pk_hex| not 32 bytes of hex.
  \item[Response 503] inbox full for this recipient
    (\verb|--max-per-recipient| reached).
\end{description}

\paragraph{GET /inbox/\{recipient\_pk\_hex\}} — recipient polls.

\begin{description}[leftmargin=2em,itemsep=0pt]
  \item[Response 200]
    \verb|{ "messages": [{ id, payload_hex, received_at_unix_ms, expires_at_unix_ms }, ...] }|
  \item[Response 400] \verb|recipient_pk_hex| malformed.
\end{description}

\paragraph{DELETE /inbox/\{recipient\_pk\_hex\}/\{message\_id\}} —
recipient marks consumed.

\begin{description}[leftmargin=2em,itemsep=0pt]
  \item[Response 200] \verb|{ "removed": bool }|
\end{description}

\paragraph{Operator endpoints.} \verb|GET /health| (uptime, deposits
total, fetches total, messages inflight) and \verb|GET /info|
(operator name, bind addr, caps, default TTL) for monitoring; no
mutation, no auth required.

\subsection{Payload Format}
\label{sec:relay-payload}

The wire-level \verb|payload_hex| is opaque to the relay. Two formats
coexist in v5.6:

\paragraph{Plain JSON (v5.3, deprecated for production).} The wallet
JSON-encodes
\[
  \{\,\texttt{coin\_secret},\,\texttt{coin\_pubkey},\,\texttt{coin\_signature},\,\texttt{denom},\,\texttt{memo}\,\}
\]
and hex-encodes the result. A relay operator with read access to the
inbox can recover all coin material. Used only when the recipient is
unable to consume sealed-box payloads (e.g. lightweight read-only
audit clients).

\paragraph{Sealed box (v5.6, default for production).} The payload is
\[
  \texttt{sealed\_blob} \;=\; \texttt{epk}\;(32)\;\|\;\texttt{ChaCha20\text{-}Poly1305(plaintext)}
\]
where the plaintext is the same JSON as above, the AEAD key is derived
via the NaCl \verb|crypto_box_seal| construction adapted for ed25519
recipient keys:

\begin{enumerate}[leftmargin=2em,itemsep=0pt]
  \item Sender samples ephemeral X25519 keypair $(esk, epk)$.
  \item Recipient's ed25519 pubkey $\to$ Montgomery form via RFC 7748
        §5 (the Edwards-to-Montgomery birational map).
  \item Shared secret $ss = esk \cdot \mathsf{recipient}_{X25519}$.
  \item $(K, N) =
        \mathsf{HKDF\text{-}SHA\text{-}256}(\,\mathsf{salt}=\texttt{"TARDUS-sealed-box-v1"},\,
                                            \mathsf{ikm}=ss,\,
                                            \mathsf{info}=epk \,\|\, \mathsf{recipient}_{X25519}\,)$.
  \item $\mathsf{ciphertext} =
        \mathsf{ChaCha20\text{-}Poly1305}(K, N, \mathsf{plaintext})$.
\end{enumerate}

Recipient inverts steps 1–4 using its own ed25519 secret key (also
canonical scalar, converted to X25519 scalar) and then runs the
matching $\mathsf{ChaCha20\text{-}Poly1305\text{-}Decrypt}$.

\subsection{Storage Backends}
\label{sec:relay-storage}

The relay's \verb|InboxStore| dispatches to one of two backends, chosen
at boot via the \verb|--storage-path| CLI flag.

\paragraph{In-memory backend (default).} \verb|HashMap<RecipientPk, Vec<StoredMessage>>|
behind a \verb|std::sync::Mutex|. Messages are lost on restart. Suited
to ephemeral testing and short-lived federation peers.

\paragraph{SQLite backend (production).} Bundled SQLite (no system
dependency), schema:

\begin{verbatim}
CREATE TABLE messages (
    id                  TEXT PRIMARY KEY,
    recipient           BLOB NOT NULL,
    payload_hex         TEXT NOT NULL,
    received_at_unix_ms INTEGER NOT NULL,
    expires_at_unix_ms  INTEGER NOT NULL
);
CREATE INDEX idx_recipient ON messages(recipient);
CREATE INDEX idx_expires   ON messages(expires_at_unix_ms);
\end{verbatim}

The two indices give $O(\log n)$ recipient-bucket scans and $O(\log n)$
TTL-prune scans even at $10^6+$ inflight messages.

\paragraph{TTL pruning.} A Tokio background task runs every $60\,\text{s}$
and issues
\verb|DELETE FROM messages WHERE expires_at_unix_ms <= now|
on SQLite or retains-by-deadline on the in-memory backend. Prune is
side-effect-free: no event log emission, no recipient notification.

\subsection{Transport Security}
\label{sec:relay-tls}

Production relays MUST run with TLS enabled. The daemon supports HTTPS
via \verb|axum-server| + \verb|rustls 0.23| (ring backend); operator
provides \verb|--tls-cert| and \verb|--tls-key| at boot. A relay with
neither flag set serves plain HTTP and refuses to accept either flag in
isolation (fail-fast misconfiguration check). v5.7 will add mTLS
support for recipient cert-based authentication on GET / DELETE.

\subsection{Security Argument}
\label{sec:relay-security}

\begin{theorem}[Sealed-box payload confidentiality]
\label{thm:relay-confidentiality}
Under assumption A1 (ECDLP) and the IND-CCA security of
ChaCha20-Poly1305, an adversary controlling the relay operator and
observing all inbox contents learns no information about the plaintext
coin material beyond the bit-length of each ciphertext, except with
probability $\leq \mathsf{Adv}^{\mathsf{IND\text{-}CCA}}_{\mathsf{ChaCha20\text{-}Poly1305}}(\lambda)
+ \mathsf{Adv}^{\mathsf{ECDLP}}_{X25519}(\lambda) + \negl(\lambda)$.
\end{theorem}

\paragraph{Reduction sketch.} The reduction is the standard
\verb|crypto_box_seal| argument: the X25519 ECDH shared secret is
indistinguishable from random under DDH (which reduces to ECDLP for
the Curve25519 family), HKDF-SHA-256 is a PRF under ROM, and
ChaCha20-Poly1305 is IND-CCA secure. Composition is via standard
hybrid arguments.

\paragraph{What the relay still learns.} The relay sees, for each
deposit: the recipient pubkey (in the URL path), the deposit timestamp,
the message TTL, and the ciphertext length. Traffic-analysis resistance
against these side channels is a deployment concern (Tor/I2P circuit,
fixed-size payload padding) and out of scope for v5.

\paragraph{What an inbox-leak does \emph{not} compromise.}
Critically, a relay compromise does not give the adversary:
(a) the recipient's ed25519 secret key (held only by the recipient's
wallet), (b) the sender's identity (no auth on POST), (c) the validator
committee's joint secret (off-chain, held by validator daemons), or
(d) any prior payments (no logging of forwarded message contents).

\subsection{Failure Mode Mapping}
\label{sec:relay-failure-modes}

Cross-reference with \Cref{sec:failure-modes}: the relay layer's
operational failure modes are extensions of the wallet/operator F-modes,
not new categories.

\begin{center}
\small
\begin{tabular}{@{}lll@{}}
\toprule
ID & Relay-layer scenario & Existing F-mode mapping \\
\midrule
R1 & Relay process crash       & F1-class (restart with persistent DB) \\
R2 & Relay disk full           & F1-class extended (SQLite write fails → 5xx) \\
R3 & Relay operator inspection & R-3 below (sealed-box mitigation) \\
R4 & Relay log seizure         & R-3 same                                \\
R5 & Inbox exhaustion attack   & F13-class (per-recipient cap, fee market) \\
R6 & Recipient secret leaked   & F8-class (wallet theft) extended         \\
R7 & Sender IP fingerprinting  & not modelled in v5 (Tor circuit at deploy) \\
\bottomrule
\end{tabular}
\end{center}

\paragraph{R-3 — relay operator payload inspection.} Mitigated by
the v5.6 sealed-box payload format. A compromised or hostile relay
operator can still observe recipient pubkeys and message metadata
(deposit timing, ciphertext size, TTL choice), but cannot recover
any coin secret or signature.

\subsection{Federation (Forward Notice)}
\label{sec:relay-federation}

A single relay is a single point of failure for delivery. The v5.7
federation extension will define a gossip protocol for relay-to-relay
forwarding so a recipient can poll any peer relay and recover messages
deposited at any federated relay. This is intentionally deferred until
the single-relay deployment pattern is stable in production.
